%!TEX program = xelatex
%!TEX root = ../all_short.tex

\chapter{Loss functions \& training: A Summary}

\section*{Loss function}

\begin{itemize}
    \item The loss function $\mathcal{L}:\mathcal{Y}\times\mathcal{Y}\mapsto\Real$ quantifies the \alert{cost} of predicting $y$ instead of the true target $t$.
    \item For an input $\x$ with true output $t$, the prediction quality is measured by:
    \begin{myequation}
    \mathcal{R}(\x, t)=L(h(\x), t)
    \end{myequation}
    \item The \alert{empirical risk} aggregates loss over the training set:
    \begin{myequation}
    \mathcal{\overline R}_{\mathcal{T}}(h)=\frac{1}{|\mathcal{T}|}\sum_{(x,t)\in\mathcal{T}}L(h(x),t)
    \end{myequation}
    \item During training, we work with parametric functions $h=h_{\theta}$, where $\Vtheta = (\theta_0, \theta_1, \ldots, \theta_d)$.
    \item Training minimizes the overall loss:
    \begin{myequation}
    \mathcal{L}(\Vtheta;\mathcal{T})=\sum_{i=1}^nL_i(\Vtheta)
    \end{myequation}
    where $L_i=L(\Vtheta; \x_i,t_i)$ for each data point.
\end{itemize}

\subsection*{Loss function minimization approaches}

\subsubsection*{Global Minimum}
\begin{itemize}
    \item A \alert{global minimum} is a point where the loss is lower than at any other point in parameter space.
    \item At any minimum, the gradient must be zero:
    \begin{myequation}
    \nabla\mathcal{L}(\Vtheta;\mathcal{T})=\Zero
    \end{myequation}
    \item The gradient $\nabla f$ is the vector of partial derivatives:
    \begin{myequation}
    \nabla f=\vcol{\frac{\partial f}{\partial x_1}}{\frac{\partial f}{\partial x_m}}
    \end{myequation}
    \item Setting gradient to zero means solving:
    \begin{myequation}
    \frac{\partial}{\partial\theta_i}\mathcal{L}(\Vtheta;\mathcal{T})=0\hspace{.5cm}\forall i
    \end{myequation}
    \item \textbf{Challenges:}
    \begin{itemize}
        \item System has multiple solutions (minima, maxima, saddle points).
        \item Analytical solutions often impossible for complex models.
    \end{itemize}
\end{itemize}

\subsection*{Gradient descent}

\begin{itemize}
    \item Iterative numerical method to find a local minimum.
    \item \textbf{Principle:} Move in the direction opposite to the gradient (steepest descent).
    \item Start from initial parameters $\Vtheta^{(0)}$ with initial error $\mathcal{\overline R}_{\mathcal{T}}(\Vtheta^{(0)})$.
    \item \textbf{Update rule:}
    \begin{myequation}
    \theta_j^{(i)} := \theta_j^{(i-1)} - \eta \frac{\partial}{\partial \theta_j }\mathcal{\overline R}_{\mathcal{T}}(\Vtheta)\Big|_{\Vtheta^{(i-1)}} =
    \theta_j^{(i-1)} - \frac{\eta}{|\mathcal{T}|}\sum_{(\x, t)\in\mathcal{T}} \frac{\partial}{\partial \theta_j} L(h_{\Vtheta}(\x), t)\Big|_{\Vtheta^{(i-1)}},
    \end{myequation}
    \item In matrix form:
    \begin{myequation}
    \Vtheta^{(i)} := \Vtheta^{(i-1)} - \eta \;\nabla\mathcal{\overline R}_{\mathcal{T}}(\Vtheta)\Big|_{\Vtheta^{(i-1)}} = 
    \Vtheta^{(i-1)} - \frac{\eta}{|\mathcal{T}|}\sum_{(\x,t)\in\mathcal{T}} \nabla L(h_{\Vtheta}(\x), t)\Big|_{\Vtheta^{(i-1)}}.
    \end{myequation}
    \item $\eta$ is the \alert{learning rate} controlling step size.
    \item \textbf{Key challenges:}
    \begin{itemize}
        \item Local minima vs global minimum.
        \item Saddle points.
        \item Convergence speed.
        \item Learning rate selection.
    \end{itemize}
\end{itemize}

\subsection*{Convexity}

\begin{itemize}
    \item A set $S\subset\Real^d$ is \alert{convex} if for any $\x_1,\x_2\in S$ and $\lambda\in (0,1)$:
    \begin{myequation}
    \lambda\x_1+(1-\lambda)\x_2\in S
    \end{myequation}
    \item A function $f(\x)$ is \alert{convex} if for all $\x_1,\x_2$ and $\lambda\in (0,1)$:
    \begin{myequation}
    f(\lambda\x_1+(1-\lambda)\x_2)\leq \lambda f(\x_1)+(1-\lambda) f(\x_2)
    \end{myequation}
    \item \alert{Strict convexity} requires strict inequality for $\x_1\neq\x_2$.
    \item \textbf{Key property:} For convex functions, any local minimum is also a global minimum.
    \item For strictly convex functions, there is exactly one local minimum (the global minimum).
    \item Properties:
    \begin{enumerate}
        \item Sum of (strictly) convex functions is (strictly) convex.
        \item Product of (strictly) convex function and positive constant is (strictly) convex.
    \end{enumerate}
    \item \textbf{Implications for empirical risk:}
    \begin{itemize}
        \item If each $L(\Vtheta; \x,t)$ is convex, empirical risk is convex.
        \item Any local minimum of empirical risk is global.
        \item If strictly convex, unique minimum exists.
    \end{itemize}
\end{itemize}

\section*{Some common loss functions}

\subsection*{Loss functions for regression}

\subsubsection*{Quadratic loss}
\begin{itemize}
    \item Definition: $L(y,t)=(y-t)^2$
    \item Empirical risk: $\mathcal{\overline R}_{\mathcal{T}}(h)=\frac{1}{|\mathcal{T}|}\sum_{(\x,t)\in\mathcal{T}}(h(\x)-t)^2$ (Mean Squared Error)
    \item For linear regression: $h(\x)=\w^T\x+b$, loss becomes:
    \begin{myequation}
    \mathcal{L}(\w,b;\mathcal{T})=\sum_{(\x,t)\in\mathcal{T}}(\w^T\x+b-t)^2
    \end{myequation}
    \item \textbf{Properties:}
    \begin{itemize}
        \item Strictly convex $\Rightarrow$ unique global minimum.
        \item Gradient is linear, making optimization tractable:
        \begin{myalign}
        \frac{\partial}{\partial w_i}\mathcal{L}(\w,b;\mathcal{T})=\sum_{(\x,t)\in\mathcal{T}}(\w^T\x+b-t)x_i\\
        \frac{\partial}{\partial b}\mathcal{L}(\w,b;\mathcal{T})=\sum_{(\x,t)\in\mathcal{T}}(\w^T\x+b-t)
        \end{myalign}
        \item Closed-form solution exists (normal equations).
        \item \textbf{Drawback:} Not robust to outliers (errors squared $\Rightarrow$ large errors dominate).
    \end{itemize}
\end{itemize}

\subsubsection*{Absolute loss}
\begin{itemize}
    \item Definition: $L(y,t)=|y-t|$ (L1 loss, Mean Absolute Error)
    \item \textbf{Properties:}
    \begin{itemize}
        \item Penalizes errors linearly, more robust to outliers.
        \item Convex but not strictly convex (flat gradient regions).
        \item Gradient is piecewise constant ($-1$ when $y < t$, $+1$ when $y > t$).
        \item Gradient undefined at $y=t$ (requires subgradients).
    \end{itemize}
\end{itemize}

\subsubsection*{Huber loss}
\begin{itemize}
    \item Definition (with parameter $\delta$):
    \begin{myequation}
    L(t,y)=
    \begin{cases}
    \frac{1}{2}(t-y)^2 & |t-y|\leq\delta \\
    \delta(|t-y|-\frac{\delta}{2}) & |t-y|>\delta
    \end{cases}
    \end{myequation}
    \item Combines quadratic behavior near zero (smooth gradients) with linear growth for large errors (robustness to outliers).
    \item $\delta$ controls transition point.
\end{itemize}

\subsection*{Loss functions for classification}

\subsubsection*{Classification approaches}
\begin{itemize}
    \item Two approaches:
    \begin{itemize}
        \item \textbf{Direct class prediction:} Output class label directly.
        \item \textbf{Probabilistic prediction:} Output probability distribution over classes.
    \end{itemize}
    \item Focus on binary classification with $t \in \{-1, +1\}$, prediction $y = h(\x)$ (confidence score, sign determines class).
\end{itemize}

\subsubsection*{0/1 loss}
\begin{itemize}
    \item Definition:
    \begin{myequation}
    L(t,y)=
    \begin{cases}
    1 &\text{sgn}(y)\not= t \\
    0 &\text{sgn}(y)= t
    \end{cases}
    \end{myequation}
    \item Compact form: $L(t,y) = {\large\mathds{1}}[t y< 0]$
    \item \textbf{Problems:}
    \begin{itemize}
        \item Non-convex.
        \item Non-smooth (discontinuous at $ty=0$).
        \item Zero gradients almost everywhere $\Rightarrow$ no gradient information.
        \item Minimization is NP-hard.
    \end{itemize}
\end{itemize}

\subsubsection*{Convex surrogate loss functions}
\begin{itemize}
    \item Surrogate loss substitutes for 0/1 loss during optimization.
    \item Required properties:
    \begin{itemize}
        \item \textbf{Upper bound property:} $L_{\text{surrogate}}(t,y) \geq {\large\mathds{1}}[ty < 0]$
        \item \textbf{Convexity:} Guarantees global minima.
        \item \textbf{Smoothness:} Enables gradient-based optimization.
    \end{itemize}
\end{itemize}

\subsubsection*{Perceptron loss}
\begin{itemize}
    \item Definition: $L(t,y)=\max{}(0, -ty)$
    \item \textbf{Properties:}
    \begin{itemize}
        \item Zero for correctly classified elements ($ty > 0$).
        \item For misclassified elements ($ty < 0$), loss $= -ty = |ty|$ (penalizes confident wrong predictions more).
        \item Continuous everywhere.
        \item Gradient continuous almost everywhere (except $ty=0$).
        \item Convex (not strictly convex).
        \item \textbf{Not} a strict surrogate (doesn't always upper bound 0/1 loss).
    \end{itemize}
\end{itemize}

\subsubsection*{Hinge loss}
\begin{itemize}
    \item Definition: $L(t,y)=\max(0, 1-ty)$
    \item \textbf{Key differences from perceptron loss:}
    \begin{itemize}
        \item Loss zero only when $ty \geq 1$ (requires confident correct predictions).
        \item For $0 < ty < 1$, loss $= 1-ty > 0$ (penalizes correct but weak predictions).
        \item For misclassified examples ($ty < 0$), loss $= 1-ty$ (grows linearly).
    \end{itemize}
    \item \textbf{Properties:}
    \begin{itemize}
        \item Continuous everywhere.
        \item Gradient continuous almost everywhere (except $ty=1$).
        \item Convex (not strictly convex due to flat region $ty > 1$).
        \item Proper surrogate loss (upper bounds 0/1 loss).
        \item Used in Support Vector Machines (encourages margin maximization).
    \end{itemize}
\end{itemize}

\subsubsection*{Subgradients}
\begin{itemize}
    \item Handle non-differentiability at singular points.
    \item A subgradient $\overline{\nabla}f|_x$ satisfies:
    \begin{myequation}
    f(x')\geq f(x)+\overline{\nabla}f|_x \cdot (x'-x) \quad \text{for all } x'
    \end{myequation}
    \item At non-differentiable points, multiple subgradients exist.
    \item For hinge loss at $ty=1$, valid subgradients are any slope in $[-t, 0]$ (when $t=1$) or $[0, -t]$ (when $t=-1$).
    \item Common choice: $\frac{\partial L_H}{\partial y}=0$ at $ty=1$, giving:
    \begin{myequation}
    \overline{\nabla}_y L_H=\begin{cases}
    -t\hspace{1cm}& ty<1\\
    0\hspace{1cm}&ty\geq 1
    \end{cases}
    \end{myequation}
\end{itemize}

\subsubsection*{Square loss in classification}
\begin{itemize}
    \item Definition: $L(t,y)=(1-ty)^2$
    \item \textbf{Properties:}
    \begin{itemize}
        \item Continuous and infinitely differentiable.
        \item Gradient continuous everywhere.
        \item Convex.
        \item \textbf{Drawbacks:}
        \begin{itemize}
            \item Not a surrogate loss (doesn't upper bound 0/1 loss everywhere).
            \item Extreme errors heavily penalized (quadratic).
            \item Symmetric penalization: even confident correct predictions ($ty \gg 1$) are penalized.
        \end{itemize}
    \end{itemize}
\end{itemize}

\subsubsection*{Log loss (cross entropy)}
\begin{itemize}
    \item Definition: $L(t,y)=\frac{1}{\log 2}\log (1+e^{-ty})$ (normalization factor sometimes omitted)
    \item Smoothed version of hinge loss.
    \item \textbf{Properties:}
    \begin{itemize}
        \item Continuous and infinitely differentiable everywhere.
        \item Gradient continuous at all points.
        \item Strictly convex $\Rightarrow$ unique global minimum.
        \item Proper surrogate loss (upper bounds 0/1 loss).
        \item As $ty \to -\infty$, loss grows linearly.
        \item As $ty \to +\infty$, loss approaches zero asymptotically.
    \end{itemize}
    \item Central to logistic regression and probabilistic classification.
\end{itemize}

\subsubsection*{Connection to information theory}
\begin{itemize}
    \item \textbf{Cross entropy:} $\entr{}{p,q} = -\ev{x\sim p}{\log q(x)}$
    \item \textbf{KL divergence:} $\dkl{p}{q} = \entr{}{p,q} - \entr{}{p}$ (non-negative, zero iff $p=q$)
    \item \textbf{Entropy:} $\entr{}{p} = -\ev{x\sim p}{\log p(x)}$ (minimum average bits needed)
    \item For binary classification with true labels $t \in \{0,1\}$ and predicted probabilities $y(\x) \in [0,1]$:
    \begin{myequation}
    \text{CE}(\mathcal{T})=-\frac{1}{|\mathcal{T}|}\sum_{(\x,t)\in\mathcal{T}}\Big(t\log y(\x)+(1-t)\log (1-y(\x))\Big)
    \end{myequation}
    \item For logistic regression ($y(\x)=\sigma(\w^T\x+b)=\frac{1}{1+e^{-(\w^T\x+b)}}$) and using $t\in\{-1,+1\}$:
    \begin{myequation}
    \text{CE}(\mathcal{T})=\frac{1}{|\mathcal{T}|}\sum_{(\x,t)\in\mathcal{T}}\log (1+e^{-t(\w^T\x+b)})
    \end{myequation}
    \item This matches log loss empirical risk (up to constant $\log 2$):
    \begin{myequation}
    \mathcal{\overline R}_{\mathcal{T}}(h)=\frac{1}{|\mathcal{T}|\log 2}\sum_{(\x,t)\in\mathcal{T}}\log (1+e^{-t(\w^T\x+b)})
    \end{myequation}
\end{itemize}

\subsubsection*{Exponential loss}
\begin{itemize}
    \item Definition: $L(t,y)=e^{-ty}$
    \item Used in boosting algorithms (AdaBoost).
    \item \textbf{Properties:}
    \begin{itemize}
        \item Continuous and infinitely differentiable.
        \item Gradient continuous everywhere.
        \item Strictly convex.
        \item Proper surrogate loss.
        \item Penalizes errors more aggressively than log loss (exponential growth).
        \item Very sensitive to outliers.
    \end{itemize}
\end{itemize}

\section*{Computing $h^*$}

\begin{itemize}
    \item Minimization of empirical risk over parameterized hypothesis class ($\theta\in\Theta=\Real^d$) is an unconstrained optimization problem.
    \item Necessary condition for optimality: stationary point where gradient vanishes:
    \begin{myequation}
    \frac{\partial}{\partial \theta_i}\mathcal{\overline R}_{\mathcal{T}}(h_\theta)=0
    \end{myequation}
    or in vector form:
    \begin{myequation}
    \nabla\mathcal{\overline R}_{\mathcal{T}}(h_\theta)=\Zero
    \end{myequation}
    \item Analytical solution generally difficult or infeasible $\Rightarrow$ numerical optimization techniques.
\end{itemize}

\subsection*{Gradient descent}

\begin{itemize}
    \item Minimize differentiable function $J(\theta)$ by iteratively moving opposite to gradient.
    \item \textbf{Update rule:}
    \begin{myalign}
    \theta^{(k+1)}&\color{evidmath}=\theta^{(k)}-\eta\,\nabla J\rvert_{\theta^{(k)}}
    \end{myalign}
    \item Component-wise:
    \begin{myequation}\color{evidmath}
    \theta_j^{(k+1)}=\theta_j^{(k)}-\eta\frac{\partial J(\theta)}{\partial \theta_j}\bigg\rvert_{\theta^{(k)}}
    \end{myequation}
    \item Four-step formulation:
    \begin{myalign}
    g_{j,k} &= \frac{\partial J(\theta)}{\partial \theta_j}\Big\rvert_{\theta^{(k)}}\\
    \eta_{j,k} &= \eta\\
    \Delta_{j,k} &= - \eta_{j,k}\, g_{j,k}\\
    \theta_j^{(k+1)} &= \theta_j^{(k)} + \Delta_{j,k}
    \end{myalign}
    \item Learning rate $\eta$ controls step size; critical for convergence speed and stability.
\end{itemize}

\subsubsection*{Batch gradient descent}
\begin{itemize}
    \item Objective: empirical risk over entire training set:
    \begin{myequation}
    J(\theta)= \mathcal{\overline R}_{\mathcal{T}}(h_\theta) = \frac{1}{|\mathcal{T}|}\sum_{(\x,t)\in\mathcal{T}} L(h_\theta(\x),t)
    \end{myequation}
    \item Gradient:
    \begin{myequation}\color{evidmath}
    g_{j,k} = \frac{1}{|\mathcal{T}|}\sum_{(\x,t)\in\mathcal{T}}\frac{\partial L(h_\theta(\x),t)}{\partial \theta_j} \Big\rvert_{\theta^{(k)}}
    \end{myequation}
    \item Update:
    \begin{myequation}\color{evidmath}
    \theta^{(k+1)}_j=\theta^{(k)}_j-\frac{\eta}{|\mathcal{T}|}\sum_{(\x,t)\in\mathcal{T}}\frac{\partial L(h_\theta(\x),t)}{\partial \theta_j} \Big\rvert_{\theta^{(k)}}
    \end{myequation}
    \item \textbf{Example (linear regression):}
    \begin{myalign}
    \theta_j^{(k+1)}&=\theta_j^{(k)}-\frac{\eta}{|\mathcal{T}|}\sum_{(\x,t)\in\mathcal{T}}\left(\sum_{s=1}^d\theta_s^{(k)}x_s+\theta_0^{(k)}-t\right)x_j \quad j=1,\ldots,d\\
    \theta_0^{(k+1)}&=\theta_0^{(k)}-\frac{\eta}{|\mathcal{T}|}\sum_{(\x,t)\in\mathcal{T}}\left(\sum_{s=1}^d\theta_s^{(k)}x_s+\theta_0^{(k)}-t\right)
    \end{myalign}
    \item \textbf{Properties:}
    \begin{itemize}
        \item Smooth, stable decrease of empirical risk.
        \item Computationally expensive for large datasets.
        \item Guaranteed convergence to global minimum if convex, local minimum if non-convex.
    \end{itemize}
    \item Pseudocode:
    \begin{lstlisting}[language=python]
for i in range(nb_epochs):
  params_grad = evaluate_gradient(loss_function, data, params)
  params = params - learning_rate * params_grad
    \end{lstlisting}
\end{itemize}

\subsubsection*{Stochastic gradient descent (SGD)}
\begin{itemize}
    \item Update using a single training example at a time:
    \begin{myequation}
    J(\theta)= L(h_\theta(\x),t)
    \end{myequation}
    \item Gradient:
    \begin{myequation}\color{evidmath}
    g_{j,k} = \frac{\partial L(h_\theta(\x),t)}{\partial \theta_j} \Big\rvert_{\theta^{(k)}}
    \end{myequation}
    \item Update:
    \begin{myequation}\color{evidmath}
    \theta_j^{(k+1)} = \theta_j^{(k)}-\eta\frac{\partial}{\partial \theta_j}L(h_\theta(\x),t)\bigg\rvert_{\theta^{(k)}}
    \end{myequation}
    \item \textbf{Properties:}
    \begin{itemize}
        \item Much faster updates, supports online learning.
        \item High variance in updates (noisy trajectory).
        \item Stochasticity can help escape shallow local minima or saddle points.
        \item With decreasing learning rate, converges almost surely to local minimum.
    \end{itemize}
    \item Pseudocode:
    \begin{lstlisting}[language=python]
for i in range(nb_epochs):
  np.random.shuffle(data)
  for example in data:
    params_grad = evaluate_gradient(loss_function, example, params)
    params = params - learning_rate * params_grad
    \end{lstlisting}
    \item \textbf{Linear regression example:}
    \begin{myalign}
    \theta_j^{(k+1)}&=\theta_j^{(k)}-\eta\left(\sum_{s=1}^d\theta_s^{(k)}x_{s}+\theta_0^{(k)}-t\right)x_{j}\\
    \theta_0^{(k+1)}&=\theta_0^{(k)}-\eta\left(\sum_{s=1}^d\theta_s^{(k)}x_{s}+\theta_0^{(k)}-t\right)
    \end{myalign}
\end{itemize}

\subsubsection*{Mini-batch gradient descent}
\begin{itemize}
    \item Intermediate between batch and SGD: use mini-batch $B$ of size $m$.
    \item Objective over mini-batch:
    \begin{myequation}
    J(\theta)= \frac{1}{m}\sum_{(\x,t)\in B} L(h_\theta(\x),t)
    \end{myequation}
    \item Gradient:
    \begin{myequation}\color{evidmath}
    g_{j,k} = \frac{1}{m}\sum_{(\x,t)\in B}\frac{\partial L(h_\theta(\x),t)}{\partial \theta_j} \Big\rvert_{\theta^{(k)}}
    \end{myequation}
    \item Update:
    \begin{myequation}\color{evidmath}
    \theta^{(k+1)}_j=\theta^{(k)}_j-\frac{\eta}{m}\sum_{(\x,t)\in B}\frac{\partial L(h_\theta(\x),t)}{\partial \theta_j}\LARGE\rvert_{\theta^{(k)}}
    \end{myequation}
    \item \textbf{Properties:}
    \begin{itemize}
        \item Limits computational cost per update.
        \item Reduces variance compared to SGD, more stable convergence.
        \item Computation can be parallelized on GPUs.
        \item Mini-batch size $m$ is a hyperparameter (typical values 32-256).
        \item Standard choice for training neural networks.
    \end{itemize}
    \item Pseudocode:
    \begin{lstlisting}[language=python]
for i in range(nb_epochs):
  np.random.shuffle(data)
  for batch in get_batches(data, batch_size=m):
    params_grad = evaluate_gradient(loss_function, batch, params)
    params = params - learning_rate * params_grad
    \end{lstlisting}
    \item \textbf{Linear regression example:}
    \begin{myalign}
    \theta_j^{(k+1)}&=\theta_j^{(k)}-\frac{\eta}{m}\sum_{(\x,t)\in B}\left(\sum_{s=1}^d\theta_s^{(k)}x_s+\theta_0^{(k)}-t\right)x_j\\
    \theta_0^{(k+1)}&=\theta_0^{(k)}-\frac{\eta}{m}\sum_{(\x,t)\in B}\left(\sum_{s=1}^d\theta_s^{(k)}x_s+\theta_0^{(k)}-t\right)
    \end{myalign}
\end{itemize}

\subsubsection*{Open issues}
\begin{itemize}
    \item Learning rate selection: too small $\Rightarrow$ slow convergence; too large $\Rightarrow$ instability/divergence.
    \item Learning rate schedules help but require predefined hyperparameters.
    \item Single global learning rate applied to all parameters may be suboptimal.
    \item Non-convex problems have many local minima and saddle points; saddle points surrounded by flat regions cause slow progress.
\end{itemize}

\subsubsection*{Momentum gradient descent}
\begin{itemize}
    \item Physical analogy: parameter vector $\theta$ as position of a body of unit mass.
    \item Force: $F(\theta)=-\eta\,\nabla J(\theta)$ (negative gradient as acceleration).
    \item Update governed by velocity, not directly by acceleration:
    \begin{myalign}
    \theta_j^{(k+1)}&=\theta_j^{(k)}+v_j^{(k+1)}\\
    v_j^{(k+1)}&=v_j^{(k)}-\eta_{j,k}\, g_{j,k}
    \end{myalign}
    \item With $v_j^{(0)}=0$ and constant learning rate:
    \begin{myequation}
    \theta_j^{(k+1)}=\theta_j^{(0)}-\eta\sum_{r=0}^kg_{j,r}
    \end{myequation}
    \item Practical version with dampening (friction) parameter $\gamma\in[0,1)$:
    \begin{myalign}
    v_j^{(k+1)}&=\gamma v_j^{(k)}-\eta g_{j,k}\\
    \theta_j^{(k+1)}&=\theta_j^{(k)}+v_j^{(k+1)}=\theta_j^{(0)}-\eta\sum_{r=1}^k\gamma^{k-r} g_{j,r}
    \end{myalign}
    \item \textbf{Properties:}
    \begin{itemize}
        \item Uses weighted sum of past gradients.
        \item Velocity increases along directions with consistent gradients.
        \item Oscillations dampened along directions with rapidly changing gradients.
        \item Reduces fluctuations, enables faster and more stable convergence.
    \end{itemize}
    \item \textbf{Linear regression example (mini-batch):}
    \begin{myalign}
    v_j^{(k+1)}&=\gamma v_j^{(k)}-\eta\sum_{(\x,t)\in B}\left(\sum_{s=1}^d \theta_j^{(k)}x_s+\theta_0^{(k)}-t\right)x_j\\
    v_0^{(k+1)}&=\gamma v_0^{(k)}-\eta\sum_{(\x,t)\in B}\left(\sum_{s=1}^d \theta_j^{(k)}x_s+\theta_0^{(k)}-t\right)\\
    \theta_j^{(k+1)}&=\theta_j^{(k)}+v_j^{(k+1)}
    \end{myalign}
\end{itemize}

\subsubsection*{Nesterov accelerated gradient descent (NAG)}
\begin{itemize}
    \item Uses look-ahead estimate $\tilde{\theta}^{(k+1)}\defeq\theta^{(k)}+\gamma v^{(k)}$ before computing gradient.
    \item Update equations:
    \begin{myalign}
    \tilde{\theta}_j^{(k+1)}&=\theta_j^{(k)}+\gamma v_j^{(k)}\\
    v_j^{(k+1)}&=\gamma v_j^{(k)} -\eta g_{j,k+1}\\
    \theta_j^{(k+1)}&=\tilde\theta_j^{(k)}+v_j^{(k+1)}
    \end{myalign}
    \item \textbf{Properties:}
    \begin{itemize}
        \item Gradient evaluated at approximated future position.
        \item Anticipates changes in velocity, correcting trajectory earlier.
        \item More responsive and often faster convergence than standard momentum.
    \end{itemize}
\end{itemize}

\subsubsection*{Dynamically updating the learning rate}
\begin{itemize}
    \item Fixed learning rate often suboptimal: larger steps early, smaller steps near optimum.
    \item \textbf{Learning rate schedules:}
    \begin{itemize}
        \item Step decay: reduce by factor $c$ every $K$ steps.
        \item Exponential decay: $\eta^{(k)} = \eta^{(0)} e^{-\alpha k}$.
        \item $1/t$ decay: $\eta^{(k)} = \frac{\eta^{(0)}}{1 + \alpha k}$.
    \end{itemize}
    \item Limitations: hyperparameters problem-dependent; same rate applied to all parameters.
    \item Adaptive methods assign parameter-specific learning rates based on gradient history.
\end{itemize}

\subsubsection*{Adagrad}
\begin{itemize}
    \item Parameter-specific learning rates that adapt to historical gradient magnitudes.
    \item For parameter $\theta_j$ at step $k$:
    \begin{myalign}
    g_{j,k} &= \frac{\partial J(\theta)}{\partial \theta_j}\Big\rvert_{\theta^{(k)}}\\
    G_{j,k} &= G_{j,k-1} + g_{j,k}^2 = \sum_{i=0}^{k} g_{j,i}^2\\
    \eta_{j,k} &= \frac{\eta}{\sqrt{G_{j,k} + \varepsilon}}\\
    \Delta_{j,k} &= - \eta_{j,k}g_{j,k}\\
    \theta_j^{(k+1)} &= \theta_j^{(k)} + \Delta_{j,k}
    \end{myalign}
    \item $\varepsilon$ small constant for numerical stability.
    \item \textbf{Properties:}
    \begin{itemize}
        \item Larger updates for rarely updated parameters; smaller updates for frequently updated ones.
        \item Well-suited for sparse data or heterogeneous features.
        \item Learning rate decays monotonically $\Rightarrow$ eventually becomes too small, halting learning.
    \end{itemize}
\end{itemize}

\subsubsection*{RMSprop}
\begin{itemize}
    \item Addresses Adagrad's monotonic decay by using exponentially decaying average of squared gradients.
    \item Decaying cumulative gradient:
    \begin{myalign}
    \tilde G_{j,k} &= \gamma \tilde G_{j,k-1} + (1 - \gamma) g_{j,k}^2 = (1 - \gamma) \sum_{i=0}^{k} \gamma^{k-i} g_{j,i}^2
    \end{myalign}
    \item Learning rate:
    \begin{myequation}
    \eta_{j,k} = \frac{\eta}{\sqrt{\tilde G_{j,k} + \varepsilon}}
    \end{myequation}
    \item Update:
    \begin{myalign}
    g_{j,k} &= \frac{\partial J(\theta)}{\partial \theta_j}\Big\rvert_{\theta^{(k)}}\\
    \tilde G_{j,k} &= \gamma \tilde G_{j,k-1} + (1 - \gamma) g_{j,k}^2\\
    \eta_{j,k} &= \frac{\eta}{\sqrt{\tilde G_{j,k} + \varepsilon}}\\
    \Delta_{j,k} &= - \eta_{j,k} g_{j,k}\\
    \theta_j^{(k+1)} &= \theta_j^{(k)} + \Delta_{j,k}
    \end{myalign}
    \item $\gamma$ controls decay rate (typically $\gamma \simeq 0.9$).
    \item \textbf{Properties:} Learning rate adapts mainly to recent gradient magnitudes; allows continued learning.
\end{itemize}

\subsubsection*{Adadelta}
\begin{itemize}
    \item Eliminates need for global learning rate $\eta$.
    \item Uses decaying average of past squared updates:
    \begin{myalign}
    \overline G_{j,k} &= \gamma \overline G_{j,k-1} + (1 - \gamma) \Delta_{j,k}^2 = (1 - \gamma) \sum_{i=0}^{k} \gamma^{k-i} \Delta_{j,i}^2
    \end{myalign}
    \item Update rule:
    \begin{myalign}
    g_{j,k} &= \frac{\partial J(\theta)}{\partial \theta_j}\Big\rvert_{\theta^{(k)}}\\
    \tilde G_{j,k} &= \gamma \tilde G_{j,k-1} + (1 - \gamma) g_{j,k}^2\\
    \eta_{j,k} &= - \frac{\sqrt{\overline G_{j,k-1} + \varepsilon}}{\sqrt{\tilde G_{j,k} + \varepsilon}}\\
    \Delta_{j,k} &= - \eta_{j,k} g_{j,k}\\
    \overline G_{j,k} &= \gamma \overline G_{j,k-1} + (1 - \gamma) \Delta_{j,k}^2\\
    \theta_j^{(k+1)} &= \theta_j^{(k)} + \Delta_{j,k}
    \end{myalign}
    \item $\gamma \simeq 0.95$, $\varepsilon \simeq 10^{-8}$.
    \item Ensures units of update $\Delta_{j,k}$ consistent with those of parameter $\theta_j$.
\end{itemize}

\subsubsection*{Adam (Adaptive Moment Estimation)}
\begin{itemize}
    \item Combines RMSprop with momentum (exponentially decaying sum of gradients).
    \item Momentum term:
    \begin{myequation}
    \tilde H_{j,k} = \beta \tilde H_{j,k-1} + (1 - \beta) g_{j,k} = (1-\beta)\sum_{i=0}^k\beta^{k-i}g_{j,i}
    \end{myequation}
    \item Bias correction (compensates for initialization at zero):
    \begin{myalign}
    \hat{G}_{j,k} &= \frac{\tilde G_{j,k}}{1 - \gamma^k}\\
    \hat{H}_{j,k} &= \frac{\tilde H_{j,k}}{1 - \beta^k}
    \end{myalign}
    \item Complete update:
    \begin{myalign}
    g_{j,k} &= \frac{\partial J(\theta)}{\partial \theta_j}\Big\rvert_{\theta^{(k)}}\\
    \tilde G_{j,k} &= \gamma \tilde G_{j,k-1} + (1 - \gamma) g_{j,k}^2\\
    \tilde H_{j,k} &= \beta \tilde H_{j,k-1} + (1 - \beta) g_{j,k}\\
    \hat{G}_{j,k} &= \frac{\tilde G_{j,k}}{1 - \gamma^k}\\
    \hat{H}_{j,k} &= \frac{\tilde H_{j,k}}{1 - \beta^k}\\
    \eta_{j,k} &= \frac{\eta}{\sqrt{\hat G_{j,k} + \varepsilon}}\\
    \Delta_{j,k} &= - \eta_{j,k} \hat H_{j,k}\\
    \theta_j^{(k+1)} &= \theta_j^{(k)} + \Delta_{j,k}
    \end{myalign}
    \item Default values: $\eta \simeq 0.001$, $\gamma \simeq 0.9$, $\beta \simeq 0.999$, $\varepsilon \simeq 10^{-8}$.
    \item Widely used due to robustness and good empirical performance.
\end{itemize}

\subsection*{Reassuming table}
\begin{center}
\setlength{\tabcolsep}{20pt}
\begin{tabular}{l|cc}
  Method & $G^{(k)}$ & $\eta^{(k)}$ \\\hline
  GD & $g^{(k)}$ & $\eta$ \\
  Adagrad & $g^{(k)}$ & $\displaystyle\frac{\eta}{RSS(g^2)}$ \\
  RMSprop & $g^{(k)}$ & $\displaystyle\frac{\eta}{RDSS(g^2,\gamma)}$ \\
  Adadelta & $g^{(k)}$ & $\displaystyle\frac{RDSS(\Delta^2,\gamma)}{RDSS(g^2,\gamma)}$ \\
  Adam & $RDSS(g,\beta)$ & $\displaystyle \frac{\eta}{RDSS(g^2,\gamma)}$ 
\end{tabular}
\end{center}

\begin{myalign}
\Delta^{(k)} &= -\eta^{(k)}G^{(k)}\\
\theta_{k+1} &= \theta^{(k)}+\Delta^{(k)}
\end{myalign}
where:
\begin{myalign}
RSS(f)&=\sqrt{f^{(0)}+f^{(1)}+\cdots+f^{(k-1)}+f^{(k)}+\varepsilon}\\
RDSS(f,\gamma)&=\sqrt{(1-\gamma)(\gamma^kf^{(0)}+\gamma^{k-1}f^{(1)}+\cdots+\gamma f^{(k-1)}+f^{(k)})+\varepsilon}
\end{myalign}

\section*{Second-Order Methods}

\begin{itemize}
    \item Seek points where $\nabla J(\theta) = 0$ using Newton-Raphson method on gradient.
    \item One-dimensional Newton:
    \begin{myequation}
    \theta_j^{(k+1)} = \theta_j^{(k)} - \frac{J'(\theta)}{J''(\theta)}\Big\rvert_{\theta^{(k)}}
    \end{myequation}
    \item Multi-dimensional Newton:
    \begin{myequation}
    \theta^{(k+1)} = \theta^{(k)} - \left( H(J)^{-1} \nabla J \right)\Big\rvert_{\theta^{(k)}}
    \end{myequation}
    where $H(J)_{ij} = \frac{\partial^2 J}{\partial \theta_i \partial \theta_j}$ is the \alert{Hessian matrix}.
    \item \textbf{Geometric interpretation:}
    \begin{itemize}
        \item Gradient Descent: approximates $J(\theta)$ with tangent plane, moves in steepest descent direction.
        \item Newton's method: approximates $J(\theta)$ with quadratic surface, moves to its minimum.
    \end{itemize}
    \item With damping factor $\eta<1$:
    \begin{myalign}
    \Delta^{(k)} &= - \left( H(J)^{-1} \nabla J \right)\bigg\rvert_{\theta^{(k)}} \\
    \theta^{(k+1)} &= \theta^{(k)} + \eta \Delta^{(k)}
    \end{myalign}
    \item \textbf{Computational drawbacks:}
    \begin{itemize}
        \item Memory: storing Hessian requires $O(n^2)$ space (prohibitive for high dimensions).
        \item Computation: inverting Hessian or solving $H \Delta = -\nabla J$ scales as $O(n^3)$.
    \end{itemize}
    \item \textbf{Quasi-Newton methods} (BFGS, L-BFGS) approximate inverse Hessian using gradient information.
    \begin{itemize}
        \item BFGS: $O(n^2)$ computation, $O(n^2)$ memory.
        \item L-BFGS: $O(mn)$ computation, $O(mn)$ memory with $m \ll n$ (linear complexity).
        \item Achieve superlinear convergence without explicit second derivatives.
    \end{itemize}
\end{itemize}

\subsection*{Optimization in deep learning}

\subsubsection*{Why first-order methods dominate}
\begin{itemize}
    \item \textbf{Scale:} Modern neural networks have billions of parameters; Hessian would have $n^2$ entries (impossible to store or invert).
    \item \textbf{Memory efficiency:} Adam maintains only two decaying sums per parameter ($O(n)$ memory).
    \item \textbf{Noise robustness:} Deep learning uses mini-batches; second-order methods are sensitive to noise in gradient estimates. Adam's decaying sums act as low-pass filters, smoothing noisy gradients.
    \item \textbf{Saddle points:} In high-dimensional spaces, saddle points are more common than local minima.
    \begin{itemize}
        \item Newton's method targets zero-gradient points, can be attracted to saddle points.
        \item Gradient descent naturally moves away from saddle points by following downward slopes.
        \item Decomposing Hessian into eigenvectors $q_i$ and eigenvalues $\lambda_i$, Newton update in direction $i$ is $\Delta \theta_i = -\frac{1}{\lambda_i} g_i$.
        \item With negative curvature ($\lambda_i < 0$), Newton moves \textit{with} the gradient, climbing toward the saddle point.
    \end{itemize}
\end{itemize}
